\chapter{Materiais e métodos}
\label{capitulo:materiais_metodos}

Este capítulo descreve a metodologia aplicada no desenvolvimento do simulador \textit{Royal Legacy}. Serão apresentados os materiais tecnológicos, abrangendo \textit{hardware} e \textit{software}, e os métodos adotados para a modelagem, levantamento de requisitos e execução das etapas do projeto. Ressalta-se que a fundamentação teórica sobre o funcionamento destas tecnologias já foi abordada na seção \ref{secao:fundamentacao_teorica}.

\section{Materiais}
\label{secao:materiais}

Os materiais empregados nesta pesquisa dividem-se em recursos físicos (\textit{hardware}), utilizados para o desenvolvimento e medição de desempenho, e lógicos (\textit{software}), que compõem o ecossistema de programação e integração do sistema.

\subsection{Recursos de \textit{Hardware}}

O desenvolvimento da aplicação, bem como os testes de estresse envolvendo o cálculo de jogadas pela Inteligência Artificial, foram conduzidos em um computador (nomeado na rede local como SERVIDOR) com as seguintes especificações:

\begin{itemize}
	\item \textbf{Processador:} AMD Ryzen 5 4600G com Radeon Graphics (3.70 GHz), arquitetura baseada em x64, contando com 6 núcleos e 12 \textit{threads}. Essa configuração foi essencial para suportar a execução da \textit{Game Engine} na via principal (\textit{main thread}) em paralelo com o motor de xadrez em segundo plano, sem causar gargalos de processamento;
	\item \textbf{Memória RAM:} 16,0 GB de memória instalada (sendo 15,4 GB utilizáveis), garantindo a alocação necessária para a manipulação de matrizes do tabuleiro e o armazenamento da árvore de busca profunda do Stockfish;
	\item \textbf{Sistema Operacional e Arquitetura:} Sistema operacional de 64 bits, adequado para compilar o executável da Godot Engine e lidar com o alto volume de instruções da Inteligência Artificial;
	\item \textbf{Armazenamento:} Unidade de armazenamento local em disco rígido/SSD, utilizada para reduzir o tempo de carregamento dos \textit{assets} visuais (temas e peças) e acelerar as chamadas de leitura do \textit{script} Python.
\end{itemize}

\subsection{Recursos de \textit{Software}}

As tecnologias e ferramentas de \textit{software} foram aplicadas de forma modularizada, separando a interface gráfica da lógica computacional:

\begin{itemize}
	\item \textbf{Godot Engine e GDScript:} O motor gráfico foi empregado como o Ambiente de Desenvolvimento Integrado (IDE) principal para a interface. O GDScript foi a linguagem utilizada para programar as regras de movimentação, o gerenciamento de cenas (\texttt{estruturador.tscn}) e o controle de áudio (\texttt{tela\_sons.gd});
	\item \textbf{Python:} Utilizado exclusivamente para o desenvolvimento do \textit{middleware} (\texttt{chess\_bridge.py}). O Python operou a biblioteca \texttt{subprocess} para abrir fluxos de leitura e escrita com a Inteligência Artificial;
	\item \textbf{Stockfish:} Empregado como o executável (\textit{backend}) que recebe os estados do tabuleiro e devolve os cálculos matemáticos das melhores jogadas durante as partidas contra a máquina;
	\item \textbf{Visual Studio Code:} Utilizado como IDE auxiliar para a escrita, formatação e depuração do código Python, facilitando a visualização dos retornos do terminal;
	\item \textbf{Protocolos FEN e UCI:} O padrão FEN foi utilizado como estrutura de dados para trafegar o posicionamento das peças entre a Godot e o Python. O protocolo UCI funcionou como a Interface de Programação de Aplicações (API) padrão para enviar comandos de nível de dificuldade ao Stockfish.
\end{itemize}

\section{Métodos}
\label{secao:metodos}

O projeto foi conduzido utilizando a metodologia de desenvolvimento iterativo e incremental. Essa escolha metodológica permitiu que o simulador de xadrez fosse construído em pequenos ciclos funcionais, garantindo que a base do jogo estivesse sólida antes da introdução de sistemas mais complexos, como a Inteligência Artificial.

\subsection{Levantamento de Requisitos e Modelagem}

Como o projeto trata-se da simulação de um jogo de tabuleiro clássico, o levantamento de requisitos não ocorreu por meio de entrevistas tradicionais com usuários, mas sim através de pesquisa documental e análise de sistemas correlatos. Os requisitos de regras lógicas foram extraídos do manual oficial da Federação Internacional de Xadrez (FIDE). Já os requisitos arquiteturais basearam-se na documentação do protocolo UCI para a troca de mensagens de texto.

Para a modelagem do sistema, optou-se pela utilização de Diagramas de Fluxo de Dados (DFD) e fluxogramas lógicos, que mapearam o caminho da informação desde o clique do usuário na interface da Godot, passando pela tradução no \textit{script} Python, até o processamento no binário do Stockfish e o seu retorno para a tela.

\subsection{Etapas de Desenvolvimento}

O desenvolvimento do \textit{Royal Legacy} foi dividido em quatro iterações principais, executadas sequencialmente:

\begin{enumerate}
	\item \textbf{Iteração 1 - Estruturação Visual:} Criação do tabuleiro em grade (\textit{grid}), instanciamento das peças 2D e implementação da lógica básica de seleção e movimentação com o \textit{mouse};
	\item \textbf{Iteração 2 - Lógica de Regras:} Codificação das validações geométricas no GDScript (movimentos permitidos por peça) e implementação de regras especiais (\textit{en passant}, roque e promoção);
	\item \textbf{Iteração 3 - Integração com a IA:} Construção do \textit{script} ponte em Python, formatação do tabuleiro em FEN e estabelecimento da comunicação assíncrona com o Stockfish sem travar a interface visual;
	\item \textbf{Iteração 4 - Polimento e Gerenciamento de Estado:} Criação do \texttt{estruturador.tscn} para transitar entre os menus, aplicação do Dicionário de dados para a troca dinâmica de temas visuais e ajustes finais no controlador de áudio.
\end{enumerate}